\begin{proof}[Proof of \cref{obj:lem:cancellation-information-envelope}]
\begin{enumerate}
\item Set
\[
b=b_m(\sigma)=\min\{1,\sigma\sqrt m/8\},
\qquad
h=h_m(\sigma)=\frac{b}{m^2},
\qquad
\lambda_{b,\sigma}=\frac{b^2}{4\sigma^2},
\qquad
N_m=(2m+1)^2-m^2.
\]
Since \(m\ge2\) and \(\sigma>0\), the support and resolution satisfy
\[
0<b\le1,\qquad h>0.
\]
Moreover, \(b\le\sigma\sqrt m/8\), so
\[
0<\lambda_{b,\sigma}
=\frac{b^2}{4\sigma^2}
\le
\frac{(\sigma\sqrt m/8)^2}{4\sigma^2}
=\frac{m}{256}.
\]
Thus the parameter conditions of \cref{obj:lem:full-marked-likelihood} hold for the alternatives constructed using \cref{obj:lem:classical-legendre-facts}. Its two likelihood bounds give
\[
\chi^2\!\left((P^+_{b,m})_{\mathrm{obs}},
             (P^-_{b,m})_{\mathrm{obs}}\right)
\le
16\kappa^2h^{2\beta}\frac{b}{N_m}
\exp\!\left(\frac{\lambda_{b,\sigma}}6\right)
\sum_{j=m-1}^{\infty}\frac{\lambda_{b,\sigma}^{\,j}}{j!}.
\]
% lean: CausalSmith.Stat.RdTruesideNoiseFrontier.noiseSupport_mem
% lean: CausalSmith.Stat.RdTruesideNoiseFrontier.noiseSupport_likelihoodLambda_le
% lean: CausalSmith.Stat.RdTruesideNoiseFrontier.full_marked_likelihood

\item Define the first retained factorial index by
\[
J_m:=m-1\in\mathbb N,
\qquad J_m\ge1,\qquad J_m\ge m/2.
\]
For every integer \(k\ge0\),
\[
\frac{\lambda_{b,\sigma}^{\,J_m+k+1}}{(J_m+k+1)!}
=
\frac{\lambda_{b,\sigma}^{\,J_m+k}}{(J_m+k)!}
\frac{\lambda_{b,\sigma}}{J_m+k+1},
\qquad
0\le\frac{\lambda_{b,\sigma}}{J_m+k+1}
\le\frac12,
\]
because \(J_m+1=m\) and \(\lambda_{b,\sigma}\le m/256\le m/2\).
Induction on \(k\), followed by summation of the geometric series, therefore yields
\[
\frac{\lambda_{b,\sigma}^{\,J_m+k}}{(J_m+k)!}
\le
\frac{\lambda_{b,\sigma}^{\,J_m}}{J_m!}\,2^{-k},
\qquad
\sum_{j=J_m}^{\infty}\frac{\lambda_{b,\sigma}^{\,j}}{j!}
\le
2\frac{\lambda_{b,\sigma}^{\,J_m}}{J_m!}.
\]
% lean: CausalSmith.Stat.RdTruesideNoiseFrontier.factorial_tail_term_le_geometric
% lean: CausalSmith.Stat.RdTruesideNoiseFrontier.noiseSupport_likelihoodTail_le

For every integer \(j\ge1\), the inequality
\(\log t\le t-1\), applied at \(t=(j+1)/j\), gives
\[
j\{\log(j+1)-\log j\}\le1.
\]
Starting from \(\log(1!)=0\ge-1\), induction gives
\[
\begin{aligned}
\log((j+1)!)
&=\log(j!)+\log(j+1)\\
&\ge j\log j-j+\log(j+1)\\
&\ge (j+1)\log(j+1)-(j+1).
\end{aligned}
\]
Consequently,
\[
\log(J_m!)\ge J_m\log J_m-J_m,
\]
and taking logarithms of the positive factorial term gives
\[
\frac{\lambda_{b,\sigma}^{\,J_m}}{J_m!}
\le
\exp\!\left[
-J_m\log\!\left(\frac{J_m}{\mathrm e\,\lambda_{b,\sigma}}\right)
\right].
\]
% lean: CausalSmith.Stat.RdTruesideNoiseFrontier.log_factorial_lower_bound
% lean: CausalSmith.Stat.RdTruesideNoiseFrontier.factorial_term_le_log_exp

To calibrate this exponent, use
\[
\frac{m}{\lambda_{b,\sigma}}\ge256,
\qquad
\log\!\left(\frac{m}{\lambda_{b,\sigma}}\right)\ge8\log2,
\qquad
\log2\ge1-\frac12=\frac12.
\]
Monotonicity of the logarithm and \(J_m\ge m/2\) imply
\[
\log\!\left(\frac{J_m}{\mathrm e\,\lambda_{b,\sigma}}\right)
\ge
\log\!\left(\frac{m}{\lambda_{b,\sigma}}\right)-\log2-1
\ge7\log2-1\ge0.
\]
Hence
\[
\begin{aligned}
&J_m\log\!\left(\frac{J_m}{\mathrm e\,\lambda_{b,\sigma}}\right)
-\frac{\lambda_{b,\sigma}}6\\
&\qquad\ge
\frac m2\left\{
\log\!\left(\frac{m}{\lambda_{b,\sigma}}\right)-\log2-1
\right\}-\frac m{1536}\\
&\qquad\ge
\frac m8\left\{1+\log\!\left(\frac{m}{\lambda_{b,\sigma}}\right)\right\}.
\end{aligned}
\]
For the last inequality, the difference between its two sides, divided by \(m\), is
\[
\begin{aligned}
&\frac38\log\!\left(\frac{m}{\lambda_{b,\sigma}}\right)
-\frac12\log2-\frac58-\frac1{1536}\\
&\qquad\ge
\frac52\log2-\frac58-\frac1{1536}
\ge\frac58-\frac1{1536}>0.
\end{aligned}
\]
Combining the preceding estimates gives
\[
\begin{aligned}
\exp\!\left(\frac{\lambda_{b,\sigma}}6\right)
\sum_{j=m-1}^{\infty}\frac{\lambda_{b,\sigma}^{\,j}}{j!}
&\le
2\exp\!\left[
\frac{\lambda_{b,\sigma}}6
-J_m\log\!\left(\frac{J_m}{\mathrm e\,\lambda_{b,\sigma}}\right)
\right]\\
&\le
2\exp\!\left[
-\frac m8\left\{1+\log\!\left(\frac{m}{\lambda_{b,\sigma}}\right)\right\}
\right].
\end{aligned}
\]
% lean: CausalSmith.Stat.RdTruesideNoiseFrontier.factorial_exponent_lower_bound
% lean: CausalSmith.Stat.RdTruesideNoiseFrontier.noiseSupport_likelihoodTail_exp_bound

\item The normalization satisfies
\[
N_m=3m^2+4m+1\ge3m^2,
\qquad
\frac{b}{N_m}\le\frac{b}{3m^2}=\frac h3.
\]
Substituting this and the factorial-tail estimate into the likelihood bound gives
\[
\begin{aligned}
\chi^2\!\left((P^+_{b,m})_{\mathrm{obs}},
             (P^-_{b,m})_{\mathrm{obs}}\right)
&\le
16\kappa^2h^{2\beta}\frac h3\,
2\exp\!\left[
-\frac m8\left\{1+\log\!\left(\frac{m}{\lambda_{b,\sigma}}\right)\right\}
\right]\\
&=
\frac{32}{3}\kappa^2h^{2\beta+1}
\exp\!\left[
-\frac m8\left\{1+\log\!\left(\frac{m}{\lambda_{b,\sigma}}\right)\right\}
\right].
\end{aligned}
\]
Since the prefactor is nonnegative and the exponential is increasing, it suffices to establish
\[
1+\Psi(h,\sigma)
\le
4m\left\{1+\log\!\left(\frac{m}{\lambda_{b,\sigma}}\right)\right\}.
\]
We split according to whether \(\sigma\sqrt m/8\ge1\).
% lean: CausalSmith.Stat.RdTruesideNoiseFrontier.noiseSupport_div_blockNorm_le
% lean: CausalSmith.Stat.RdTruesideNoiseFrontier.cancellation_information_degree_bound
% lean: CausalSmith.Stat.RdTruesideNoiseFrontier.cancellation_information_envelope

\item Suppose first that \(1\le\sigma\sqrt m/8\). Then
\[
b=1,\qquad
\sigma^2m\ge64,\qquad
h=\frac1{m^2},\qquad
\lambda_{b,\sigma}=\frac1{4\sigma^2}.
\]
Squaring \(\sigma^2m\ge64\) shows that
\[
\sigma^4m^2\ge4096>1,
\qquad
h<\sigma^4\le\sigma,
\]
where the final inequality follows from \(0<\sigma\le1\).
Thus the compact branch of the cost applies. Using
\[
\sqrt h=\frac1m,\qquad h^{-1/2}=m,
\]
we obtain
\[
\begin{aligned}
1+\Psi(h,\sigma)
&=m\{1+\log(\sigma^2m)\}\\
&\le m\{1+\log(4\sigma^2m)\}\\
&=m\left\{1+\log\!\left(\frac{m}{\lambda_{b,\sigma}}\right)\right\}.
\end{aligned}
\]
The logarithmic comparison follows from
\(0<\sigma^2m\le4\sigma^2m\).
Also, \(\lambda_{b,\sigma}\le m/256\) implies
\[
0\le
m\left\{1+\log\!\left(\frac{m}{\lambda_{b,\sigma}}\right)\right\}
\le
4m\left\{1+\log\!\left(\frac{m}{\lambda_{b,\sigma}}\right)\right\}.
\]
This proves the required cost comparison in the saturated case, including the equality \(\sigma\sqrt m/8=1\).
% lean: CausalSmith.Stat.RdTruesideNoiseFrontier.noiseSupport_saturated_threshold
% lean: CausalSmith.Stat.RdTruesideNoiseFrontier.noiseCost_saturated_eq
% lean: CausalSmith.Stat.RdTruesideNoiseFrontier.noiseCost_saturated_le_degree

\item Suppose now that \(\sigma\sqrt m/8<1\). Direct substitution in the support and likelihood scales gives
\[
b=\frac{\sigma\sqrt m}{8},
\qquad
\sigma^2m=64b^2,
\qquad
\lambda_{b,\sigma}=\frac m{256},
\qquad
\frac{m}{\lambda_{b,\sigma}}=256.
\]
We follow the three branches of the cost definition.

If \(\sigma\le h\), then
\[
1+\Psi(h,\sigma)=1
\le4m(1+\log256),
\]
because \(m\ge2\) and \(\log256\ge0\).
% lean: CausalSmith.Stat.RdTruesideNoiseFrontier.noiseSupport_unsaturated_identities

If \(\sigma>h\) and \(\sigma^4\le h\), then
\[
1+\Psi(h,\sigma)=(\sigma/h)^{2/3}.
\]
Indeed, \(m\ge2\) gives \(\sqrt m\le m<8m^2\), and hence
\[
h=\frac{\sigma\sqrt m}{8m^2}<\sigma.
\]
The scale identity \(\sigma^2m=64b^2\) yields
\[
\left\{(\sigma/h)^{2/3}\right\}^{3}
=(\sigma/h)^2
=\frac{\sigma^2m^4}{b^2}
=64m^3=(4m)^3.
\]
Both quantities whose cubes are equal are positive, so
\[
1+\Psi(h,\sigma)=4m\le4m(1+\log256).
\]
% lean: CausalSmith.Stat.RdTruesideNoiseFrontier.noiseCost_unsaturated_intermediate_eq

Finally, suppose that \(\sigma>h\) and \(h<\sigma^4\). The scale identity gives
\[
\sigma^4m^2=4096b^4,
\qquad
b<\sigma^4m^2=4096b^4.
\]
These inequalities force \(b\ge1/16\): if \(b<1/16\), then
\[
b^3\le\frac1{4096},
\qquad
4096b^4\le b,
\]
contradicting the strict inequality above. Therefore
\[
\sqrt h=\frac{\sqrt b}{m}\ge\frac1{4m},
\qquad
h^{-1/2}=\frac1{\sqrt h}\le4m.
\]
On the other hand, \(0<b\le1\) implies
\[
b^3\le1,\qquad b^4\le b,\qquad b^2\le\sqrt b.
\]
Consequently, the logarithm's argument satisfies
\[
0<
\frac{\sigma^2}{\sqrt h}
=\frac{\sigma^2m}{\sqrt b}
=\frac{64b^2}{\sqrt b}
\le64\le256.
\]
Monotonicity of the logarithm and \(1+\log256\ge0\) now give
\[
\begin{aligned}
1+\Psi(h,\sigma)
&=h^{-1/2}\left\{1+\log\!\left(\frac{\sigma^2}{\sqrt h}\right)\right\}\\
&\le h^{-1/2}(1+\log256)\\
&\le4m(1+\log256).
\end{aligned}
\]
Thus every unsaturated branch satisfies
\[
1+\Psi(h,\sigma)
\le
4m\left\{1+\log\!\left(\frac{m}{\lambda_{b,\sigma}}\right)\right\}.
\]
The weak inequalities in the first two branches include their boundary values.
% lean: CausalSmith.Stat.RdTruesideNoiseFrontier.noiseCost_compact_support_bound
% lean: CausalSmith.Stat.RdTruesideNoiseFrontier.noiseCost_unsaturated_le_degree

\item The cost comparison in both support cases implies
\[
-\frac m8\left\{1+\log\!\left(\frac{m}{\lambda_{b,\sigma}}\right)\right\}
\le-\frac{1+\Psi(h,\sigma)}{32}.
\]
Applying monotonicity of the exponential to the degree bound therefore proves
\[
\chi^2\!\left((P^+_{b,m})_{\mathrm{obs}},
             (P^-_{b,m})_{\mathrm{obs}}\right)
\le
\frac{32}{3}\kappa^2h^{2\beta+1}
\exp\!\left(-\frac{1+\Psi(h,\sigma)}{32}\right).
\]
% lean: CausalSmith.Stat.RdTruesideNoiseFrontier.cancellation_information_envelope
\end{enumerate}
\end{proof}
